\bookchapter{Capitolo 9: Ufficio mobile e apprendimento}

\chapterintro{Anche io un tempo pensavo che lavorare sul telefono fosse solo “arrangiarsi”, finché una volta in viaggio d'affari ho dimenticato il computer e ho dovuto modificare un bando e firmare un contratto direttamente sul telefono. A dire il vero, il collo di bottiglia del lavoro su telefono spesso non sta nel telefono stesso, ma nel fatto che non si conoscono i suoi limiti di capacità.}

\booksection{Elaborazione documenti: editare e creare sul palmo della mano}

Modificare documenti sullo schermo del telefono è diverso rispetto al computer. Il problema maggiore è lo schermo piccolo, le dita grosse, l'impaginazione che slitta facilmente. Ma alcuni accorgimenti migliorano molto l'esperienza:

\begin{itemize}
  \item \textbf{WPS Office}: la localizzazione è la migliore, pienamente compatibile con i formati Microsoft (docx/xlsx/pptx), con un ricco set di strumenti PDF integrati. Mi piace molto la sua funzione "Note vocali" - quando si esamina un contratto e non è comodo digitare, basta dettare alcune opinioni al telefono, e lui le converte automaticamente in annotazioni testuali inserite nel documento. Inoltre WPS integra una gran quantità di modelli gratuiti, fa risparmiare molto tempo per curriculum e PPT di reportistica.
  \item \textbf{Microsoft 365}: se la tua azienda usa Office originale, questa è la scelta migliore. La compatibilità del formato con la versione desktop è quasi perfetta, non si verificano quei leggeri spostamenti di layout complesso che a volte capitano con WPS. Profondamente integrato con OneDrive: il documento che hai modificato a metà sul computer, sul telefono si apre già nella versione più recente.
  \item \textbf{Shimo Docs/Tencent Docs}: queste due sono rappresentative della "collaborazione online". Adatte a più persone che modificano lo stesso documento contemporaneamente, con le modifiche di ciascuno mostrate in tempo reale. Ho usato Shimo Docs per preparare una guida di viaggio con un amico, entrambi inserivamo informazioni nella tabella contemporaneamente, con un'efficienza di gran lunga superiore rispetto al farlo da soli.
\end{itemize}

\textbf{Flusso di lavoro pratico}: la mia abitudine attuale è - i file importanti vanno su OneDrive, l'editing quotidiano con WPS, la collaborazione di squadra su Shimo Docs. Da quando i file sono in cloud, non si è più verificato l'imbarazzo di “aver dimenticato la chiavetta USB”.

\booksection{App di note: dagli appunti rapidi al sistema di conoscenza}

Per quanto riguarda prendere appunti, molti vogliono subito costruire un sistema complesso, ma poi non resistono oltre tre giorni. Il mio consiglio è prima semplice poi complesso: prima costruisci l'abitudine di “annotare”, poi pensa al problema della “gestione”.

\begin{itemize}
  \item \textbf{Appunti rapidi - il sistema nativo è sufficiente}: le Note di iOS si avviano rapidissime e supportano scansione documenti, scrittura a mano, registrazione audio. Le note dei vari marchi Android non sono da meno, possono essere avviate rapidamente anche a schermo bloccato. Questi strumenti sono adatti per annotare idee temporanee, liste della spesa, punti chiave delle riunioni. La chiave è l'avvio rapido: viene in mente qualcosa e la si annota, senza aspettare il caricamento dell'app.
  \item \textbf{Leggi dopo - raccolta di informazioni}: la funzione "Finestra mobile" di WeChat è in realtà uno strumento nascosto. Vedi un buon articolo senza tempo per leggerlo, clicca in alto a destra "Finestra mobile", si ridurrà a un puntino attaccato al bordo dello schermo, riapribile in qualsiasi momento. Molto più comodo che inoltrarlo all'assistente di trasferimento file. Inoltre, strumenti come Cubox o Pocket permettono di raccogliere su più piattaforme i buoni contenuti visti ovunque.
  \item \textbf{Costruire un sistema - Notion e Feishu}: quando le note sparse diventano difficili da trovare, puoi passare a Notion. È davvero flessibile: puoi creare un "database di note di lettura" organizzato automaticamente per titolo del libro, autore, valutazione; o un "tabellone di progetto" per visualizzare l'avanzamento dei compiti di lavoro. Anche le "Tabelle multidimensionali" di Feishu sono simili e in Cina l'accesso è più veloce. Uso Notion per gestire tutti i miei materiali di scrittura, note di studio e avanzamento dei progetti da oltre tre anni.
\end{itemize}

\textbf{Una strategia pratica per le note}: la mia abitudine è una “tris” - idee temporanee nelle note di sistema (avvio immediato), articoli da leggere nella finestra mobile di WeChat (a portata di mano), conoscenza da accumulare a lungo termine su Notion (sistematica). Non cercare fin da subito di infilare tutto in un solo strumento: è un modo facile per mollare.

\booksection{Scansione e OCR: portare il mondo cartaceo nel telefono}

La scansione con il telefono, se usata bene, può davvero risparmiare molti fastidi di andare in copisteria. L'ultima volta che ho dovuto presentare i documenti per il visto, li ho sistemati tutti con il telefono.

\begin{itemize}
  \item \textbf{Utenti iPhone}: la funzione di scansione integrata nelle Note di iOS è già molto buona. Apri Note tocca l'icona fotocamera scegli "Scansiona documento", puntalo sul file e riconosce automaticamente i bordi correggendo la prospettiva. Dopo lo scatto puoi regolare i colori - in modalità "Bianco e nero" il risultato è quello più vicino a uno scanner, "Scala di grigi" è adatto a documenti con immagini.
  \item \textbf{Utenti Android}: molte fotocamere dei marchi hanno una modalità "Scansione documento" integrata. Se non la trovi, scarica Microsoft Lens dallo store (gratis e integrato con OneDrive).
  \item \textbf{Riconoscimento testo OCR}: per i documenti scansionati, se serve estrarre il testo per modificarlo, serve l'OCR. La funzione "Estrai testo" integrata in WeChat è la più semplice - premi a lungo sull'immagine e si usa, la precisione del riconoscimento è buona. Per grandi quantità, puoi usare l'app Baimiao o CamScanner. Per il cinese, il motore OCR di Baidu è buono; anche la conversione immagine-testo di WPS è precisa.
  \item \textbf{Piccoli trucchi di scansione}: scatta in un luogo ben illuminato, possibilmente con una lampada da tavolo laterale per evitare l'ombra della mano sul documento. Se il foglio è piegato agli angoli, appiattiscilo con la mano prima di scattare. Per documenti di più pagine puoi scansionare in sequenza e unire in un unico PDF, senza salvarli pagina per pagina.
\end{itemize}

\booksection{Quale cloud scegliere: non c'è il migliore, solo il più adatto a te}

Sulla scelta del cloud ho riflettuto a lungo, alla fine ho capito che dipende dal tuo “ecosistema principale”.

\begin{itemize}
  \item \textbf{Ecosistema Apple iCloud}: sincronizzazione così fluida da non accorgersene. Scatti una foto sul telefono e dopo pochi minuti è già sul computer. E Apple è la più rigorosa nella protezione della privacy. Ma gli svantaggi sono evidenti: la condivisione multipiattaforma con utenti Android o Windows è scomoda, i 5GB gratuiti sono pochi, espandere costa.
  \item \textbf{Windows principale OneDrive}: se usi un computer Windows + Office 365, OneDrive è una coppia naturale. Salvi un documento in Word e finisce automaticamente in cloud. 1TB insieme all'abbonamento Office è molto conveniente. Anche il backup automatico delle foto dall'app mobile è molto comodo.
  \item \textbf{Condivisione nazionale Baidu Netdisk}: a dire il vero, il limite di velocità di download di Baidu Netdisk è stato spesso criticato, ma il suo ecosistema di condivisione è davvero imbattibile. Gran parte dei materiali di studio e delle risorse video sul Web cinese sono condivisi tramite Baidu Netdisk. Se scarichi solo occasionalmente file condivisi da altri, si può usare anche senza abbonamento. Se sei un raccoglitore incallito di risorse, conviene attivare l'abbonamento: il tempo risparmiato vale di più.
  \item \textbf{Alibaba Cloud Drive/Quark Drive}: sono le novità degli ultimi due anni. Alibaba Cloud Drive non limita la velocità di upload/download (al momento), Quark Drive è strettamente integrato con il browser Quark, puoi salvare direttamente nel cloud durante la ricerca di materiali, molto comodo.
\end{itemize}

\textbf{La mia soluzione combinata}: iCloud per foto personali e file quotidiani (privacy prima), Baidu Netdisk per risorse e ricezione di condivisioni (ecosistema prima), OneDrive per documenti di lavoro (integrazione Office prima). Tre cloud, ciascuno con il suo compito, senza interferenze.

\booksection{Uno scenario reale di lavoro mobile}

Ti racconto un'esperienza personale. Una volta in viaggio d'affari in un'altra città, il cliente ha chiesto improvvisamente di modificare le clausole del contratto. Il mio flusso di lavoro è stato:

\begin{enumerate}
  \item In hotel ho ricevuto via WeChat le osservazioni del cliente (messaggio testuale + foto di annotazioni a mano).
  \item Aperto OneDrive per recuperare il contratto originale, aperto e modificato con WPS, seguendo i messaggi WeChat del cliente per modificare le clausole una per una.
  \item Dopo le modifiche, esportato in PDF con la funzione "Esporta come PDF" di WPS e inviato al gruppo WeChat del cliente per conferma.
  \item Il cliente ha dato una seconda tornata di osservazioni, con lo strumento penna di WPS ho cerchiato le posizioni da modificare direttamente sul PDF, con note vocali per spiegare le idee di modifica.
  \item Dopo conferma, ho usato la funzione "Firma elettronica" di WPS per firmare online e inviato all'altra parte per controfirma.
  \item Infine ho archiviato il PDF del contratto firmato sia su Baidu Netdisk che su iCloud, e contemporaneamente ho annotato una riga nel diario di lavoro su Notion.
\end{enumerate}

L'intero processo è durato circa quaranta minuti, senza mai accendere il computer. A dire il vero, tre anni fa non avrei mai creduto che il telefono potesse fare tutto questo, ma oggi è davvero possibile. La chiave è familiarizzare in anticipo con questi strumenti, per non doverli scaricare e registrare in modo disordinato quando servono.

\begin{tcbpractice}{Prova: costruisci il tuo flusso di lavoro mobile}
\begin{enumerate}
  \item Scegli la tua app di documenti principale (WPS o Microsoft 365), accedi e collega il cloud.
  \item Apri Note o Memo, usa la funzione di scansione per scansionare un documento, esportalo in PDF e invialo a te stesso.
  \item Scegli un'app di cloud, impostala per il backup automatico delle foto del telefono (scegli "Backup solo sotto Wi-Fi" per risparmiare dati).
  \item Crea una nuova pagina in Notion o Feishu e prova a impostare un semplice modello "Diario di lavoro".
\end{enumerate}

\end{tcbpractice}
